% from ega3-1-fr.tex
\oldpage[0\textsubscript{III}]{20}
topologie $\mathfrak n$-adique, et on sait que cette dernière est \emph{séparée}
(\hyperref[0.7.3.5-fr]{$0_{\mathrm I}$, 7.3.5}). D'ailleurs, comme $M/M_i$ est
$A$-plat, on a
\[
  f^i(M\otimes_A N)=\operatorname{Im}(M_i\otimes_A N\to M\otimes_A N)
  =\operatorname{Ker}(M\otimes_A N\to(M/M_i)\otimes_A N)
\]
(\hyperref[0.6.1.2-fr]{$0_{\mathrm I}$, 6.1.2}). Soient alors $N$ un
$A$-module de type fini, $N'$ un sous-module de $N$, $j:N'\to N$ l'injection
canonique~; dans le diagramme commutatif
\[
\xymatrix{
  M\otimes_A N'\ar[r]\ar[d]_{1_M\otimes j} &
  (M/M_i)\otimes_A N'\ar[d]^{1_{M/M_i}\otimes j}\\
  M\otimes_A N\ar[r] & (M/M_i)\otimes_A N
}
\]
$1_{M/M_i}\otimes j$ est injectif puisque $M/M_i$ est $A$-plat~; on en
conclut que
\[
  \operatorname{Ker}(M\otimes_A N'\to M\otimes_A N)
  \subset\operatorname{Ker}(M\otimes_A N'\to(M/M_i)\otimes_A N')
\]
quel que soit $i$~; puisque l'intersection des seconds membres est réduite à
$0$ comme on l'a vu plus haut, il en est de même du premier membre, et par
suite $M$ est $A$-plat.

\begin{proposition}[10.2.8]
\label{0.10.2.8-fr}
Soient $A$ un anneau noethérien réduit, $M$ un $A$-module de type fini. On
suppose que pour toute $A$-algèbre $B$, qui est un anneau de valuation
discrète, $M\otimes_A B$ soit un $B$-module plat (donc libre
(\hyperref[0.10.1.3-fr]{10.1.3})). Alors $M$ est un $A$-module plat.
\end{proposition}

On sait que pour que $M$ soit plat, il faut et il suffit que pour tout idéal
maximal $\mathfrak m$ de $A$, $M_{\mathfrak m}$ soit un
$A_{\mathfrak m}$-module plat
(\hyperref[0.6.3.3-fr]{$0_{\mathrm I}$, 6.3.3})~; on peut donc se borner au cas
où $A$ est \emph{local} (\hyperref[0.1.2.8-fr]{$0_{\mathrm I}$, 1.2.8}).
Soient alors $\mathfrak m$ l'idéal maximal de $A$, $\mathfrak p_i$
$(1\leq i\leq r)$ ses idéaux premiers minimaux, $k$ le corps résiduel
$A/\mathfrak m$. On sait (\hyperref[II.7.1.7-fr]{II, 7.1.7}) qu'il existe
pour chaque $i$ un anneau de valuation discrète $B_i$, ayant même corps des
fractions $K_i$ que l'anneau intègre $A/\mathfrak p_i$, et dominant ce
dernier. Posons $M_i=M\otimes_A B_i$. Par hypothèse, $M_i$ est libre sur
$B_i$, donc on a, en désignant par $k_i$ le corps résiduel de $B_i$
\[
\label{0.10.2.8.1-fr}
  \operatorname{rg}_{k_i}(M_i\otimes_{B_i}k_i)
  =\operatorname{rg}_{K_i}(M_i\otimes_{B_i}K_i).
  \tag{10.2.8.1}
\]
Mais il est clair que l'homomorphisme composé
$A\to A/\mathfrak p_i\to B_i$ est local, donc $k$ est une extension de $k_i$,
et l'on a
$M_i\otimes_{B_i}k_i=M\otimes_A k_i=(M\otimes_A k)\otimes_k k_i$, et par
ailleurs $M_i\otimes_{B_i}K_i=M\otimes_A K_i$. L'égalité
(\hyperref[0.10.2.8.1-fr]{10.2.8.1}) entraîne donc
\[
  \operatorname{rg}_k(M\otimes_A k)=\operatorname{rg}_{K_i}(M\otimes_A K_i)
  \qquad\text{pour }1\leq i\leq r
\]
et comme $A$ est réduit, on sait que cette condition entraîne que $M$ est un
$A$-module \emph{libre} (Bourbaki, \emph{Alg. comm.}, chap. II, \S~3,
n\textsuperscript{o}~2, prop. 7).

\subsection{Existence d'extensions plates d'anneaux locaux}
\label{subsection:0.10.3-fr}

\begin{proposition}[10.3.1]
\label{0.10.3.1-fr}
Soient $A$ un anneau local noethérien, $\mathfrak J$ son idéal maximal,
$k=A/\mathfrak J$ son corps résiduel. Soit $K$ une extension du corps $k$. Il
existe un homomorphisme local de $A$ dans un anneau local noethérien $B$, tel
que $B/\mathfrak J B$ soit $k$-isomorphe à $K$, et que $B$ soit un
$A$-module plat.
\end{proposition}

Nous démontrerons cette proposition en plusieurs étapes.

\begin{env}[10.3.1.1]
\label{0.10.3.1.1-fr}
Supposons d'abord que $K=k(T)$, où $T$ est une indéterminée. Dans l'anneau de
polynômes $A'=A[T]$, considérons l'idéal premier
$\mathfrak J'=\mathfrak J A'$ formé des
